\chapter{Conclusion} \label{cap:conclusion}

\section{Final Remarks} \label{sec:final-remarks}

This work built a deterministic Forex evaluation pipeline that combines price-derived and news-derived features, then used it to test whether adding a news family to a price-only strategy improves cost-adjusted return. The executed study is a fixed-model EUR/USD evaluation: one baseline and one hybrid strategy replayed on LEAN under a common execution model, followed by registered follow-ups over one trading year (Chapter~\ref{cap:experimental-evaluation}). The ten-year, two-pair, multi-source design with rolling retraining and combinatorial cross-validation remains a target and not a result; it is mentioned in Section~\ref{subsec:future-unexecuted}.

\subsection{Answer to the research question} \label{subsec:answer-rq}

In the configuration tested, the hybrid strategy did not outperform the baseline, and no strategy variant showed an edge. Both March--October 2016 simulations on frozen models lose about 88\% of the deposit, and the news family changes the outcome by about \$65 over 1{,}029 trades. Fifty-one further tuning variants on the H4, H1 and M15 clocks confirm this: none is profitable in November 2016, the one month that no earlier decision had read, and the loss grows with the number of trades. Hypothesis~H1 is therefore not supported. The result is bounded by its scope: one currency pair, one year, a daily event-intensity index rather than article text, and a single chronological split per experiment. It does not show that news has no value for Forex prediction in general.

\subsection{Scientific findings} \label{subsec:scientific-findings}

The experiments yield six findings about method and signal, independent of the software that produced them.

\begin{enumerate}
    \item \textbf{A positive return does not imply predictive ability.} The single positive cell, the reversed-sign news rule at $+12.2\,\%$, is the euro's decline with leverage: its always-short control is statistically the same account ($p = 0.92$), and its long-side mirror loses $48.5\,\%$. Without a drift control this cell would have read as a discovery.
    \item \textbf{Thresholds calibrated on one month do not survive a level shift.} The event intensity never returned to the January-2016 level on which its low cut was calibrated, so one sign was long-only and the other short-only in effect.
    \item \textbf{The learned models carried no information.} The F7 models are at or above the coin-flip log-loss of $0.693$ from March onward; the later losses come from the money-management template acting on uninformative predictions, not from a model going stale.
    \item \textbf{A signal must be traded on its own horizon.} The signal bars of the news rule carry a tilt of about four pips over four hours, which the reported trades never captured because the exit template holds positions for days; the signal then becomes the drift.
    \item \textbf{Filters and candlestick patterns do not survive multiple-comparison control.} The one nominally positive pattern of a 22-pattern screen is indistinguishable from the chance outcome of screening 22 candidates (Section~\ref{subsec:candles-multiple-comparisons}), and the six TA-Lib patterns read as continuation signals lean the wrong way at H1.
    \item \textbf{A uniformly negative result has two explanations that this design cannot separate:} no exploitable signal at the tested horizon, or a method unable to find it. The candidate causes of the second are specific: hand-picked thresholds never fitted on held-out data, a veto-based chain in which any one filter can kill a trade instead of letting weak signals offset one another, and a one-year, single-pair window.
\end{enumerate}

\subsection{Infrastructure and reproducibility contributions} \label{subsec:infra-contributions}

Separately from the findings above, the work leaves a reusable and auditable experimental apparatus: a reproducible data-to-replay pipeline; a deterministic filter chain with a per-decision audit trail; frozen model artifacts identified by hash; byte-for-byte replays (twelve models retrained from scratch are identical to the reported ones outside their provenance, fifteen cells reproduce byte for byte, and the pre-registered primary test returns the same $p = 0.149$); a repository registry that labels each job registered, exploratory or ad hoc; a training-to-engine feature-parity test; and corrected statistical tooling for stationary-bootstrap and deflated-Sharpe analysis (Section~\ref{subsec:rule-significance}). These contributions allow the findings to be checked; they are not themselves evidence about the market.

\subsection{Limitations} \label{subsec:limitations}

The meta-learner, as implemented, receives only the probability outputs of the LightGBM family models, so its inputs share a common probability scale; differences in scale concern the raw features inside the family models and any future early-fusion combiner. Repeated threshold, gate and strategy selection on the same months prevents treating the reused windows as confirmatory out-of-sample evidence. Template-derived M1 and H4 sweeps do not reproduce the author's Spring strategy, which is unpublished. The adaptive-retraining evidence remains incomplete: the reported 60-day monthly refit did not beat the frozen comparator at either tested clock, and the remaining registered policies have not yet produced a full comparison.

\section{Future Work} \label{sec:future-work}

\subsection{Registered follow-ups} \label{subsec:future-registered}

Four follow-ups were registered as planned studies on September 28, 2026, each with controls and endpoints fixed before any run. Two have since produced results (Chapter~\ref{cap:experimental-evaluation}): the frozen policy of the adaptive-training study over the registered year ($-97.8\,\%$, driven by a 26\,\% win rate while the payoff ratio stayed favorable at 1.97) and the candlestick-catalogue screen, whose single-chain run with all 22 patterns lost $99.54\,\%$ over 2015 (M1) and $32.79\,\%$ over the registered one-year window (H1). Two remain open: a confluence chain whose sequence is fixed in advance (momentum context, event-intensity trigger on trailing quantiles, optional indicator confirmation, risk guard, exit on the signal's own horizon), and a second exogenous feature family built from other markets after \citeonline{laidi2009intermarket}. The remaining four adaptive-training policies (refreshed thresholds, rolling, expanding, exponentially weighted) are also outstanding, so no ranking across policies is yet possible. The Fibonacci confluence field was implemented but never enabled in a reported backtest; it is a gap in the evaluation and not a null result.

\subsection{Components designed but not executed} \label{subsec:future-unexecuted}

The original design contained components that no reported experiment exercised. They are not part of the evidence and are listed here only as directions:

\begin{itemize}
    \item \textbf{Text sentiment and more news sources:} article-level sentiment (a lexicon baseline, a FinBERT-class scorer, later an LLM agent), the GPR index as a model input, and further corpora. This is the most direct test of whether article text, as opposed to coded events, adds information.
    \item \textbf{A wider filter catalogue:} additional confirmation, regime and consensus gates, and a weighted terminal rule that lets weak signals offset each other, which the veto chain does not allow.
    \item \textbf{A fuller validation protocol:} combinatorial purged cross-validation with embargo, rolling refits over a multi-year window, and a studentized bootstrap for short windows.
    \item \textbf{More instruments and longer windows:} USD/JPY and other pairs, a currency-strength family, and other asset classes through adapter substitution.
\end{itemize}
